% Candidate rules for order structure and status.  Written by hand; see docs/one-order-coverage.md.

\section{Order structure and status}

\begin{definition}[Order structure (C-01)]
  \label{def:c01-structure}
  \lean{Rules.OneOrder.Service, Rules.OneOrder.ServiceStatus, Rules.OneOrder.OrderItem, Rules.OneOrder.Order}
  \uses{def:order, def:order-item, def:passenger}
  An order holds one or more order items, and each order item holds one or
  more services.  A service is for one passenger on one flight segment, and is
  the smallest part of an order that can be delivered on its own.  Each
  service has a service status, such as \texttt{CONFIRMED} or
  \texttt{TRANSFERRED}, and a delivery status, such as
  \texttt{READY TO PROCEED}.

  Source: \href{https://standards.atlassian.net/wiki/spaces/EASD/pages/574849062/Order}{EASD guide, Order}.
\end{definition}

\begin{definition}[Offer items for passengers and segments]
  \label{def:one-offer}
  \lean{Rules.OneOrder.OfferItem, Rules.OneOrder.Offer}
  \uses{def:offer, def:passenger, def:money, def:instant}
  An offer item is priced for one or more passengers on one or more flight
  segments.  An offer holds its items, an expiry instant and a currency.

  Source: \href{https://standards.atlassian.net/wiki/spaces/EASD/pages/574849062/Order}{EASD guide, Order}.
\end{definition}

\begin{definition}[Creating an order from an offer]
  \label{def:one-accept}
  \lean{Rules.OneOrder.accept}
  \uses{def:one-offer, def:c01-structure, def:rejection}
  Accepting an offer before it expires creates an order with one order item
  for each offer item, and one confirmed service for each passenger and
  segment of that item.  The checks before creating the order are the same as
  in the pilot rule: expiry, then an empty offer, then mixed currencies.

  Source: \href{https://standards.atlassian.net/wiki/spaces/EASD/pages/574849062/Order}{EASD guide, Order}.
\end{definition}

\begin{theorem}[One order item per offer item (C-02)]
  \label{thm:c02-item-per-offer-item}
  \lean{Rules.OneOrder.accept_item_per_offer_item}
  \uses{def:one-accept, def:c01-structure}
  Accepting an offer creates exactly one order item for each offer item.

  Source: \href{https://standards.atlassian.net/wiki/spaces/EASD/pages/574849062/Order}{EASD guide, Order}.
\end{theorem}
\begin{proof}
  \uses{def:one-accept}
  The order's items are the offer's items, each turned into an order item that
  records where it came from.
\end{proof}

\begin{theorem}[One service per passenger and segment]
  \label{thm:one-services-per-passenger-segment}
  \lean{Rules.OneOrder.accept_services_per_passenger_segment}
  \uses{def:one-accept, def:c01-structure}
  Each order item created from an offer item holds one service for each
  passenger and flight segment of that offer item.

  Source: \href{https://standards.atlassian.net/wiki/spaces/EASD/pages/574849062/Order}{EASD guide, Order}.
\end{theorem}
\begin{proof}
  \uses{def:one-accept}
  The services are built from the pairs of a passenger and a segment.
\end{proof}

\begin{theorem}[Order item IDs are unique (C-03)]
  \label{thm:c03-unique-item-ids}
  \lean{Rules.OneOrder.accept_item_ids_unique}
  \uses{def:one-accept, def:c01-structure}
  Accepting an offer whose item IDs are distinct gives the order's items
  distinct IDs.  Whether every change keeps them distinct belongs with the
  change rules (C-18).

  Source: \href{https://standards.atlassian.net/wiki/spaces/EASD/pages/574849062/Order}{EASD guide, Order}.
\end{theorem}
\begin{proof}
  \uses{def:one-accept}
  Each order item ID is the order ID, a hyphen and the offer item ID, so
  distinct offer item IDs give distinct order item IDs.
\end{proof}

\begin{theorem}[Order items are not split (C-04)]
  \label{thm:c04-item-unsplit}
  \uses{def:c01-structure, def:c18-change}
  The price is held on the order item, and no operation changes part of an
  item.  A change replaces whole items.

  Source: \href{https://standards.atlassian.net/wiki/spaces/EASD/pages/574849062/Order}{EASD guide, Order}.
\end{theorem}

\begin{definition}[Services that hold value (C-05)]
  \label{def:c05-holds-value}
  \lean{Rules.OneOrder.ValueReading}
  \uses{def:c08-delivery-status}
  Whether a service still holds value, decided by its delivery status.  The
  source settles five statuses: \texttt{CONFIRMED}, \texttt{READY TO PROCEED}
  and \texttt{REMOVED HOLD VALUE} hold value, and \texttt{DELIVERED} and
  \texttt{REMOVED} do not.  The model takes the answer for the other eight as a
  parameter, a reading of the guide, and a guarantee about value counts as
  proven only when it holds for every reading.

  Open point: which of the other eight delivery statuses hold value, for
  example \texttt{EXPIRED}, \texttt{NOT CLAIMED} and \texttt{SUSPENDED}.

  Source: \href{https://standards.atlassian.net/wiki/spaces/EASD/pages/574849062/Order}{EASD guide, Order}.
\end{definition}

\begin{definition}[Order item status (C-06)]
  \label{def:c06-item-status}
  \lean{Rules.OneOrder.itemStatuses}
  \uses{def:c01-structure, def:c05-holds-value}
  An order item is \texttt{ACTIVE} while at least one of its services holds
  value.  It is \texttt{CANCELLED} once all its services are cancelled or
  transferred and none holds value.  It is \texttt{CLOSED} once at least one
  service is delivered and every other service is delivered or removed.

  Source: \href{https://standards.atlassian.net/wiki/spaces/EASD/pages/574849062/Order}{EASD guide, Order}.
\end{definition}

\begin{theorem}[Active items are those holding value]
  \label{thm:one-item-active}
  \lean{Rules.OneOrder.item_active_iff}
  \uses{def:c06-item-status}
  An order item is marked active exactly when at least one of its services
  holds value.

  Source: \href{https://standards.atlassian.net/wiki/spaces/EASD/pages/574849062/Order}{EASD guide, Order}.
\end{theorem}
\begin{proof}
  \uses{def:c06-item-status}
  Unfold the status rule.
\end{proof}

\begin{theorem}[Cancelled items as the guide defines them]
  \label{thm:one-item-cancelled}
  \lean{Rules.OneOrder.item_cancelled_iff}
  \uses{def:c06-item-status}
  An order item is marked cancelled exactly when every service is cancelled
  or transferred and none holds value.

  Source: \href{https://standards.atlassian.net/wiki/spaces/EASD/pages/574849062/Order}{EASD guide, Order}.
\end{theorem}
\begin{proof}
  \uses{def:c06-item-status}
  Unfold the status rule.
\end{proof}

\begin{theorem}[Closed items as the guide defines them]
  \label{thm:one-item-closed}
  \lean{Rules.OneOrder.item_closed_iff}
  \uses{def:c06-item-status}
  An order item is marked closed exactly when at least one service is
  delivered and every service is delivered or removed.

  Source: \href{https://standards.atlassian.net/wiki/spaces/EASD/pages/574849062/Order}{EASD guide, Order}.
\end{theorem}
\begin{proof}
  \uses{def:c06-item-status}
  Unfold the status rule.
\end{proof}

\begin{theorem}[Active items are never closed]
  \label{thm:one-active-not-closed}
  \lean{Rules.OneOrder.active_not_closed}
  \uses{def:c06-item-status}
  Under every reading of the guide, an active item is not closed.

  Source: \href{https://standards.atlassian.net/wiki/spaces/EASD/pages/574849062/Order}{EASD guide, Order}.
\end{theorem}
\begin{proof}
  \uses{def:c06-item-status, def:c05-holds-value}
  A closed item's services are all delivered or removed, and neither holds
  value.
\end{proof}

\begin{theorem}[Active items are never cancelled]
  \label{thm:one-active-not-cancelled}
  \lean{Rules.OneOrder.active_not_cancelled}
  \uses{def:c06-item-status}
  An active item is not cancelled.

  Source: \href{https://standards.atlassian.net/wiki/spaces/EASD/pages/574849062/Order}{EASD guide, Order}.
\end{theorem}
\begin{proof}
  \uses{def:c06-item-status}
  The cancelled status requires that no service holds value.
\end{proof}

\begin{theorem}[Each item has exactly one status]
  \label{thm:one-item-one-status}
  \lean{Rules.OneOrder.item_has_one_status}
  \uses{def:c06-item-status, thm:one-active-not-closed, thm:one-active-not-cancelled}
  Every order item with a service has exactly one of the three statuses.

  Counterexample: an item whose one service was delivered and then
  transferred meets the conditions for both \texttt{CANCELLED} and
  \texttt{CLOSED}.  An item whose one service is \texttt{SUSPENDED} has no
  status at all under a reading where a suspended service holds no value.

  Open point: which status wins when the conditions overlap, whether a
  delivered service may be transferred at all, and what status an item has
  when none of the conditions holds.

  Source: \href{https://standards.atlassian.net/wiki/spaces/EASD/pages/574849062/Order}{EASD guide, Order}.
\end{theorem}
\begin{proof}
  \uses{def:c06-item-status}
  Not proven.  The counterexample above shows the statement is false for the
  rule as written.
\end{proof}

\begin{definition}[Order status (C-07)]
  \label{def:c07-order-status}
  \lean{Rules.OneOrder.orderStatuses}
  \uses{def:c06-item-status, def:c05-holds-value}
  An order is \texttt{OPENED} while at least one of its services holds value.
  It is \texttt{CLOSED} once all its items are closed or cancelled and all its
  services are delivered or removed.

  Source: \href{https://standards.atlassian.net/wiki/spaces/EASD/pages/574849062/Order}{EASD guide, Order}.
\end{definition}

\begin{theorem}[Opened orders are those holding value]
  \label{thm:one-order-opened}
  \lean{Rules.OneOrder.order_opened_iff}
  \uses{def:c07-order-status}
  An order is marked opened exactly when at least one of its services holds
  value.

  Source: \href{https://standards.atlassian.net/wiki/spaces/EASD/pages/574849062/Order}{EASD guide, Order}.
\end{theorem}
\begin{proof}
  \uses{def:c07-order-status}
  Unfold the status rule.
\end{proof}

\begin{theorem}[Closed orders as the guide defines them]
  \label{thm:one-order-closed}
  \lean{Rules.OneOrder.order_closed_iff}
  \uses{def:c07-order-status}
  An order is marked closed exactly when every item is closed or cancelled
  and every service is delivered or removed.

  Source: \href{https://standards.atlassian.net/wiki/spaces/EASD/pages/574849062/Order}{EASD guide, Order}.
\end{theorem}
\begin{proof}
  \uses{def:c07-order-status}
  Unfold the status rule.
\end{proof}

\begin{theorem}[New orders are opened]
  \label{thm:one-new-order-opened}
  \lean{Rules.OneOrder.accept_order_opened}
  \uses{def:one-accept, def:c07-order-status}
  Under every reading of the guide, an order created from an offer whose
  items each have a passenger and a segment is opened and not closed.

  Source: \href{https://standards.atlassian.net/wiki/spaces/EASD/pages/574849062/Order}{EASD guide, Order}.
\end{theorem}
\begin{proof}
  \uses{def:one-accept, def:c05-holds-value}
  Its services start confirmed, and a confirmed service holds value.
\end{proof}

\begin{theorem}[Orders are never both opened and closed]
  \label{thm:one-order-exclusive}
  \lean{Rules.OneOrder.opened_closed_exclusive}
  \uses{def:c07-order-status}
  Under every reading of the guide, no order is both opened and closed.

  Source: \href{https://standards.atlassian.net/wiki/spaces/EASD/pages/574849062/Order}{EASD guide, Order}.
\end{theorem}
\begin{proof}
  \uses{def:c07-order-status, def:c05-holds-value}
  An opened order has a service that holds value; a closed order's services
  are all delivered or removed, and neither holds value.
\end{proof}

\begin{theorem}[Every order is opened or closed]
  \label{thm:one-order-has-status}
  \lean{Rules.OneOrder.order_opened_or_closed}
  \uses{def:c07-order-status}
  Every order with at least one service is opened or closed.

  Counterexample: under a reading where an expired service holds no value, an
  order whose only service expired is neither opened nor closed.

  Open point: whether \texttt{EXPIRED} and \texttt{UNABLE TO DELIVER} hold
  value, and whether an order with such a service should be closed.

  Source: \href{https://standards.atlassian.net/wiki/spaces/EASD/pages/574849062/Order}{EASD guide, Order}.
\end{theorem}
\begin{proof}
  \uses{def:c07-order-status}
  Not proven.  The counterexample above shows the statement is false for the
  rule as written.
\end{proof}
